% ===================================================================
%  BAGAN No. 10 — Laut. Pelabuhan.
%
%  Traduction, faite en 2026, en indonesien d'aujourd'hui, sur l'ido
%  de texto/io. Regles de la colonne : voir 10-bagan-01.tex. Une seule
%  numerotation.
%
%  LE MEME PARTAGE QU'AU TABLEAU 10 OURDOU, DANS UNE TROISIEME FAMILLE
%  DE LANGUES. Tout ce que la marine a voile connaissait porte un nom
%  malais de fond : perahu le bateau, layar la voile, tiang le mat,
%  jangkar l'ancre, geladak le pont, haluan la proue, buritan la
%  poupe, kemudi le gouvernail, dayung la rame, dermaga le quai,
%  pelabuhan le port, selat le detroit, teluk la baie, tanjung le cap,
%  semenanjung la presqu'ile, tanah genting l'isthme, pulau l'ile,
%  karang le recif, pasang surut la maree, mercusuar le phare. Tout ce
%  que la vapeur a invente vient du neerlandais ou de l'anglais :
%  kapal uap, kapal penjelajah, torpedo, ketel, dok, derek. Le
%  releve a maintenant vu ce partage en ourdou, en indonesien et — par
%  l'autre bout — au chemin de fer indonesien : il ne tient ni a la
%  famille ni a la geographie, il tient a la DATE.
%
%  MAIS L'AMIRAL (84) FAIT L'INVERSE EXACT DE L'OURDOU. La colonne
%  ourdoue ecrit ایڈمرل et note que le persan avait pourtant امیر
%  البحر, ancetre du mot « admiral » lui-meme, qu'elle refuse de
%  ressusciter parce que personne ne l'emploie. L'indonesien n'a rien
%  a ressusciter : il dit laksamana, et il le dit tous les jours —
%  c'est le grade officiel de la marine indonesienne. Le mot vient du
%  Ramayana, ou Lakshmana est le frere fidele de Rama : un titre naval
%  tire d'une epopee sanskrite, la ou l'Europe a pris le sien a
%  l'arabe. Meme fonction, deux mondes.
%
%  ET L'ARCHIPEL EXPORTE SES BATEAUX. perahu a donne « proa » a
%  l'anglais et au francais ; sampan est passe partout ; le jong
%  javanais a donne « jonque ». Avec damar au tableau 9 et beranda au
%  tableau 4, cela fait quatre mots que cette colonne rend au
%  vocabulaire europeen plutot qu'elle ne lui emprunte — et trois des
%  quatre sont des choses qui se transportent par mer.
%
%  selat (98), LE DETROIT, MERITE SA LIGNE. Le mot est celui du detroit
%  de Malacca, par ou passe aujourd'hui un quart du commerce mondial.
%  Une planche belge de 1926 met un detroit au fond du paysage comme
%  un accident du relief ; pour le lecteur indonesien, c'est la chose
%  la plus ordinaire de son pays — dix-sept mille iles n'ont pas
%  d'autre facon de se toucher.
% ===================================================================

%%K t10-apar-1 apar
\VUsaut{4.0mm}
\VUtitre{15.0pt}{60}{BAGAN No. 10}
\VUsaut{3.0mm}
\VUpk{11.4pt}{40}{Laut. --- Pelabuhan.}
\VUsaut{3.0mm}

%%K t10-01-1 p
1. --- Pada musim panas, sementara sebagian orang mencari ketenangan
dan kesejukan di gunung, yang lain lebih suka pergi ke pantai. Itulah
yang kami lakukan tahun lalu, sebab kami sangat menyukai
\VUgras{Samudra Atlantik} \textsuperscript{(1)}.

%%K t10-02-1 p
2. --- Bagi saya sendiri, saya merasakan kesenangan besar memandang
hamparan air yang luar biasa luas itu, yang terbentang sampai
\VUgras{kaki langit} \textsuperscript{(2)}, dan melihat
\VUgras{kapal-kapal uap} \textsuperscript{(3)} yang sangat besar
berangkat untuk \VUgras{pelayaran} yang panjang, tempat para penumpang
yang menuju Amerika naik.

%%K t10-03-1 p
3. --- Karena itu ayah saya, yang mengetahui kesukaan saya, selalu
memilih untuk menghabiskan liburan sebuah pantai yang sangat tenang,
yang terletak dekat salah satu \VUgras{pangkalan angkatan laut} kami
yang besar.

%%K t10-03-2 p
Selama kami tinggal di sana yang terakhir, \VUgras{armada} datang
berlabuh di balik \VUgras{tanggul laut} \textsuperscript{(4)} yang
sangat besar, yang menutup dan melindungi \VUgras{pelabuhan militer}
\textsuperscript{(5)}, dan saya berhasil mengagumi sepuas hati
bermacam-macam \VUgras{kapal perang} : \VUgras{kapal selam}
\textsuperscript{(10)} kecil yang menyelam dan lenyap di bawah air,
dan juga \VUgras{kapal berlapis baja} \textsuperscript{(9)} yang sangat
besar, yang sampai sekarang hampir hanya takut pada serangan
\VUgras{kapal torpedo} \textsuperscript{(8)}.

%%K t10-tit-2 sub
\VUsaut{3.0mm}
\VUpk{10.2pt}{40}{Kapal-kapal kecil.}
\VUsaut{3.0mm}

%%K t10-04-1 p
4. --- Kapal itu sudah pandai dengan sangat \VUgras{cekatan} menemukan
titik lemah lapisan baja yang tebal untuk menempelkan di sana
\VUgras{torpedo} yang berbahaya, yang ledakannya menenggelamkan kapal ;
namun ia kurang menakutkan daripada kapal selam, sebab kapal
penghancur torpedo dengan mudah mengejarnya.

%%K t10-05-1 p
5. --- Saya juga dapat melihat \VUgras{kapal-kapal penjelajah}
\textsuperscript{(6)} berlapis baja yang cepat, hampir sekebal kapal
berlapis baja yang sesungguhnya, beberapa \VUgras{kapal pengangkut}
\textsuperscript{(12)}, tiga atau empat \VUgras{kapal meriam}
\textsuperscript{(7)}, dan akhirnya sebuah \VUgras{fregat}
\textsuperscript{(11)}. \VUgras{Pemandangan armada itu, ketika ia
bergerak untuk mengangkat sauh, sungguh yang paling menarik.}

%%K t10-06-1 p
6. --- Alangkah berbedanya bangunan antara tumpukan baja yang besar
itu dan \VUgras{kapal tiga tiang} yang anggun atau \VUgras{kapal
layar} \textsuperscript{(13)} yang bagus, yang sampai sekarang masih
dipakai dalam armada dagang, yang saya lihat masuk ke pelabuhan dengan
layar terkembang penuh ketika saya berjalan-jalan di \VUgras{dermaga}
\textsuperscript{(14)}. Sungguh, kapal-kapal itu kehilangan banyak
keuntungannya \VUgras{ketika} \VUgras{angin} berlawanan. Mereka harus
menurunkan semua layar untuk masuk kembali ke kolam, dan sebuah
\VUgras{kapal tunda} \textsuperscript{(16)} diperlukan untuk menjemput
dan membawa mereka, seperti \VUgras{tongkang} \textsuperscript{(17)}
biasa, ke kade.

%%K t10-tit-3 sub
\VUsaut{3.0mm}
\VUpk{10.2pt}{40}{Pelabuhan dagang.}
\VUsaut{3.0mm}

%%K t10-07-1 p
7. --- Yang menarik pada pelabuhan yang saya bicarakan itu ialah bahwa
ia berisi bukan hanya pelabuhan militer, yang ditata secara buatan
dengan pekerjaan raksasa, melainkan juga sebuah \VUgras{pelabuhan
dagang} yang menakjubkan, yang terletak di \VUgras{muara} sebuah
sungai besar.

%%K t10-08-1 p
8. --- Lidah tanah yang memisahkan kedua kolam itu diperkuat dan
dijaga oleh serangkaian \VUgras{baterai meriam} dan \VUgras{benteng}
yang menjorok ke laut. Orang memerlukan izin khusus untuk berjalan di
\VUgras{tembok kota} \textsuperscript{(15)}. Saya dengan mudah
memperolehnya lewat paman saya, yang menjadi insinyur \VUgras{galangan
kapal} \textsuperscript{(18)}, dan saya pergi melihat kapal-kapal yang
sedang dibangun di bawah hanggar yang luas.

%%K t10-09-1 p
9. --- Saya bahkan hadir pada peluncuran sebuah kapal berlapis baja,
dan saya ingat saat yang mengharukan yang dirasakan seluruh kerumunan
ketika massa raksasa itu, dilepaskan pada dirinya sendiri, mulai
meluncur di atas rangka peluncurnya dan datang terapung di atas air
sungai.

%%K t10-10-1 p
10. --- Untuk sampai ke galangan, orang melintasi \VUgras{lapangan
latihan} \textsuperscript{(19)}, tempat para prajurit garnisun berlatih
setiap hari ; orang melewati sebuah \VUgras{jembatan kecil}
\textsuperscript{(20)} dari besi, yang menghubungkan lapangan parade
dengan \VUgras{gudang senjata} \textsuperscript{(21)}.

%%K t10-tit-4 sub
\VUsaut{3.0mm}
\VUpk{10.2pt}{40}{Gudang senjata dan dok.}
\VUsaut{3.0mm}

%%K t10-11-1 p
11. --- Bersama paman saya, saya mengunjungi bangunan besar itu, yang
penuh dengan \VUgras{meriam} \textsuperscript{(22)}, \VUgras{peluru
meriam} \textsuperscript{(23)} dan \VUgras{peluru howitzer.} Saya juga
melihat \VUgras{pabrik logam} \textsuperscript{(24)} tempat
\VUgras{ketel} \textsuperscript{(25)} kapal uap dibuat.

%%K t10-12-1 p
12. --- Sangat dekat dari situ ada \VUgras{dok-dok}
\textsuperscript{(26)} — dok perbaikan — dan \VUgras{dok-dok terapung}
\textsuperscript{(27)} yang sangat menarik, tempat saya dapat melihat
dari sangat dekat, pada sebuah kapal yang sedang diperbaiki,
\VUgras{baling-baling} \textsuperscript{(28)}, \VUgras{lunas}
\textsuperscript{(29)} dan \VUgras{kemudi} \textsuperscript{(30)} sebuah
perahu.

%%K t10-13-1 p
13. --- Pelabuhan dagang sama layaknya dipelajari seperti pelabuhan
militer, dan saya menghabiskan berjam-jam melongo di kade tempat
\VUgras{kapal-kapal roda} \textsuperscript{(31)} yang menyusuri sungai
ke hulu ditambatkan, sambil memandang \VUgras{derek-derek bertiang}
\textsuperscript{(32)} bekerja, yang mengangkat muatan yang sangat
besar seperti bulu.

%%K t10-tit-5 sub
\VUsaut{4.0mm}
\VUpk{10.2pt}{40}{Pandu pelabuhan.}
\VUsaut{3.0mm}

%%K t10-14-1 p
14. --- Saya mengagumi \VUgras{kapal-kapal lintas Atlantik}
\textsuperscript{(34)} yang keluar dari kolam, dengan
\VUgras{bendera-bendera} \textsuperscript{(35)} mereka terkembang di
udara, yang dipandu oleh seorang \VUgras{pandu pelabuhan}
\textsuperscript{(36)} yang cakap, yang secara menakjubkan tahu
menyusuri \VUgras{alur pelayaran} yang ditandai dengan
\VUgras{pelampung} \textsuperscript{(40)} dan \VUgras{tiang tanda,}
sambil menghindari \VUgras{tongkang-tongkang lebar}
\textsuperscript{(41)} yang dengan susah payah dikemudikan oleh
layarnya yang tunggal dan bersegi empat. Saya membedakan di
\VUgras{haluan} \textsuperscript{(37)} — bagian depan kapal —
\VUgras{kabin-kabin} \textsuperscript{(39)} kelas dua, dan di
\VUgras{buritan} \textsuperscript{(38)} — bagian belakang — ruang-ruang
duduk untuk penumpang kelas satu.

%%K t10-15-1 p
15. --- Kadang-kadang saya juga hadir pada pemuatan sebuah
\VUgras{tongkang} \textsuperscript{(42)} yang baru saja diisi
barang-barang dari \VUgras{gudang} \textsuperscript{(43)} dan
\VUgras{stasiun laut} \textsuperscript{(44)}, yang dibawa oleh
kereta-kereta yang tak terhitung banyaknya di \VUgras{rel dermaga}
\textsuperscript{(45)}.

%%K t10-tit-6 sub
\VUsaut{3.0mm}
\VUpk{10.2pt}{40}{Bersampan.}
\VUsaut{3.0mm}

%%K t10-16-1 p
16. --- Saya sering berperahu di \VUgras{kolam pelabuhan}
\textsuperscript{(53)}, tempat \VUgras{perahu-perahu}
\textsuperscript{(46)} datang berlindung, dan menjadi kegembiraan bagi
saya mengayuh \VUgras{dayung} \textsuperscript{(47)}. Saya menaiki
\VUgras{geladak} \textsuperscript{(48)} sebuah \VUgras{perahu layar}
\textsuperscript{(49)} ; saya memanjat, dengan \VUgras{tali-tali besar}
\textsuperscript{(52)}, ke \VUgras{tiang} \textsuperscript{(51)} sampai
ke \VUgras{andang-andang} \textsuperscript{(50)} ; kemudian saya
mengulangi tamasya saya \VUgras{dengan sampan} \textsuperscript{(54)} di
antara \VUgras{perahu-perahu nelayan} \textsuperscript{(55)} yang berat,
tempat \VUgras{jaring} \textsuperscript{(56)} dijemur.

%%K t10-17-1 p
17. --- Di \VUgras{kade} \textsuperscript{(58)} lewat di depan saya
\VUgras{barang-barang} \textsuperscript{(59)} yang paling
bermacam-macam, yang dikemas dalam \VUgras{peti} \textsuperscript{(60)}
atau dalam \VUgras{bal} \textsuperscript{(61)}. Semua itu merupakan
\VUgras{muatan} \textsuperscript{(63)} tongkang-tongkang besar, dan
\VUgras{derek-derek} \textsuperscript{(62)} yang perkasa mengangkat dan
meletakkannya di atas \VUgras{truk} \textsuperscript{(64)}, sementara
\VUgras{para pengangkut} melaporkannya dan membayar beanya di
\VUgras{kantor bea cukai} \textsuperscript{(65)}.

%%K t10-18-1 p
18. --- Akhirnya, ketika saya sampai di \VUgras{jembatan putar}
\textsuperscript{(66)} atau di \VUgras{pintu air} \textsuperscript{(69)},
sambil lewat di depan \VUgras{kapal keruk} \textsuperscript{(68)} yang
membersihkan \VUgras{terusan} \textsuperscript{(67)}, saya turun ke
dermaga di dekat \VUgras{tiang isyarat} \textsuperscript{(70)}.

%%K t10-19-1 p
19. --- Di sisi lain dermaga itulah \VUgras{sekoci-sekoci}
\textsuperscript{(71)} yang membawa para perwira armada ke darat
merapat. Sungguh pemandangan yang patut dikagumi melihat para pelaut
mengangkat dayung mereka dengan berirama, sementara perahu yang
diangkat \VUgras{ombak} \textsuperscript{(72)} merapat dengan mudah ke
tangga tempat mendarat. Dari tempat ini orang dapat mengamati dengan
lebih baik \VUgras{pasang surut} dan gerak \VUgras{pasang} dan
\VUgras{surut} di \VUgras{pantai} \textsuperscript{(73)}.

%%K t10-tit-7 sub
\VUsaut{3.0mm}
\VUpk{10.2pt}{40}{Gedung perkumpulan.}
\VUsaut{3.0mm}

%%K t10-20-1 p
20. --- Untuk pulang ke rumah, saya memakai \VUgras{trem listrik}
\textsuperscript{(74)} yang \VUgras{haltenya} \textsuperscript{(75)}
berada dekat \VUgras{tiang} yang ditanam di depan gedung perkumpulan
para perwira.

%%K t10-20-2 p
Dari \VUgras{balkon} \textsuperscript{(76)} gedung itu, yang dihiasi
\VUgras{jangkar} \textsuperscript{(77)}, meriam dan peluru meriam,
sering saya melihat kumpulan perwira laut yang megah :
\VUgras{letnan-letnan} \textsuperscript{(78)} dengan \VUgras{pedang}
mereka, \VUgras{kapten-kapten artileri} \textsuperscript{(79)} dan
\VUgras{infanteri} \textsuperscript{(80)} dengan \VUgras{pedang
lengkung} mereka. Di belakang mereka berdiri seorang \VUgras{kelasi}
\textsuperscript{(81)} yang membawakan \VUgras{teropong} kepada mereka
bila mereka hendak membedakan apa yang terjadi di kejauhan.

%%K t10-21-1 p
21. --- Saya juga melihat \VUgras{letnan-letnan muda}
\textsuperscript{(82)} yang menerima perintah dari
\VUgras{komandannya} \textsuperscript{(83)} atau dari
\VUgras{laksamana} \textsuperscript{(84)}, sementara seorang
\VUgras{kepala regu} \textsuperscript{(85)} menunggu untuk membawa
perintah itu ke kapal.

%%K t10-22-1 p
22. --- Dari balkon itu pemandangannya megah : orang menguasai seluruh
pantai dengan \VUgras{tendanya} \textsuperscript{(86)} dan \VUgras{bilik
gantinya} \textsuperscript{(87)}, dan pada jam mandi orang melihat
\VUgras{orang-orang yang mandi} \textsuperscript{(88)} bermain dengan
riang di depan \VUgras{kasino} \textsuperscript{(89)}.

%%K t10-23-1 p
23. --- Di ujung pantai, daratan membentuk sebuah \VUgras{semenanjung}
\textsuperscript{(91)} kecil yang \VUgras{berbatu karang,} tempat
\VUgras{gelombang} \textsuperscript{(92)} datang memecah, dan di atasnya
dibangun sebuah \VUgras{mercusuar} \textsuperscript{(93)} yang pada
malam hari menunjukkan jalan kepada para pelaut. Semenanjung itu
berhubungan dengan daratan hanya melalui sebuah \VUgras{tanah genting}
\textsuperscript{(90)} yang sangat sempit.

%%K t10-tit-8 sub
\VUsaut{3.0mm}
\VUpk{10.2pt}{40}{Pemandangan umum pesisir.}
\VUsaut{3.0mm}

%%K t10-24-1 p
24. --- \VUgras{Tebing} \textsuperscript{(94)} itu terbentang
panjang, terbelah oleh laut, yang dengan sesuka hati menggambar
padanya serangkaian \VUgras{teluk} \textsuperscript{(95)} dan
\VUgras{tanjung,} yang membuatnya sangat elok.

%%K t10-24-2 p
Di \VUgras{dataran tinggi} \textsuperscript{(96)}, yang menguasai
\VUgras{selat} \textsuperscript{(98)}, berdiri sebuah \VUgras{benteng}
\textsuperscript{(97)} yang sangat penting dan beberapa « vila ».

%%K t10-25-1 p
25. --- Selat itu memisahkan dari benua sebuah \VUgras{pulau}
\textsuperscript{(99)} yang sangat berbahaya, yang dikelilingi
\VUgras{karang} \textsuperscript{(100)} dan \VUgras{pecahan ombak,}
tempat perahu dan bahkan kapal besar kadang-kadang terdampar. Meskipun
pertolongan disiapkan dan \VUgras{sekoci penolong} datang dengan cepat,
peristiwa \VUgras{karam} itu mengerikan, dan selama \VUgras{badai}
ekuinoks beberapa kapal binasa bersama orang dan barangnya.

%%K t10-25-2 p
Meskipun bertetangga dengan bahaya itu, \VUgras{pelabuhan} itu sangat
sering disinggahi para pelaut.

\VUsaut{6.0mm}
\VUfilet{22mm}
